% =====================================================================
\section{Последовательный порт: разговор с компьютером}
% =====================================================================
\lessonintro
  {понять, как устроен UART, научиться выводить отладочные сообщения и принимать команды с компьютера}
  {урок 1: неблокирующий цикл, светодиод лампы}
  {лампа включается и выключается командами \code{on} / \code{off} из Serial Monitor}
  {то же, что в уроке 1; для последнего раздела --- второе устройство с UART (по желанию)}

Прежде чем добавлять в лампу кнопки и датчики, нужен способ видеть, что происходит внутри
программы. Этим способом будет последовательный порт: через него мы будем проверять
результат в каждом следующем уроке.

\subsection{Кадр UART}

UART --- асинхронный протокол: у приёмника и передатчика нет общего тактового сигнала, поэтому обе стороны
должны заранее договориться о скорости (бодах). Линия в покое находится в состоянии HIGH.

\begin{figure}[H]
\centering
\begin{tikzpicture}[x=1.05cm,y=1cm]
  \def\hi{1}
  \draw[very thick,accent] (0,\hi) -- (1,\hi) -- (1,0) -- (2,0) -- (2,\hi) -- (3,\hi) -- (3,0) -- (8,0) -- (8,\hi) -- (9,\hi) -- (9,0) -- (10,0) -- (10,\hi) -- (12,\hi);
  \foreach \x in {1,...,11} \draw[gray!50,dashed] (\x,-0.3) -- (\x,1.4);
  \foreach \x/\l in {1.5/Старт,2.5/D0,3.5/D1,4.5/D2,5.5/D3,6.5/D4,7.5/D5,8.5/D6,9.5/D7,10.5/Стоп}
    \node[font=\sffamily\scriptsize] at (\x,1.65) {\l};
  \foreach \x/\b in {2.5/1,3.5/0,4.5/0,5.5/0,6.5/0,7.5/0,8.5/1,9.5/0}
    \node[font=\ttfamily\small,text=esp] at (\x,-0.5) {\b};
  \node[font=\sffamily\scriptsize] at (0.5,1.65) {покой};
  \node[font=\sffamily\scriptsize] at (11.5,1.65) {покой};
  \draw[decorate,decoration={brace,mirror,amplitude=4pt}] (2,-0.8) -- (10,-0.8)
    node[midway,below=5pt,font=\sffamily\small]{8 бит данных, младший бит первым: \texttt{0x41} = 'A'};
  \draw[<->] (1,-1.8) -- node[below,font=\scriptsize]{$t_\text{бит} = 1/115200 \approx \SI{8.68}{\micro s}$} (2,-1.8);
\end{tikzpicture}
\caption{Передача символа «A» в формате 8N1 (8 бит данных, без контроля чётности, 1 стоп-бит)}
\end{figure}

\subsection{Serial и Serial Monitor}

У ESP32-S3 три аппаратных UART, а в Arduino им соответствуют объекты \code{Serial}:

\begin{center}\small
\begin{tabular}{@{}llll@{}}
\toprule
Объект Arduino & Контроллер & Выводы по умолчанию & Назначение \\
\midrule
\code{Serial} & USB CDC или UART0 & USB / 43, 44 & связь с компьютером \\
\code{Serial0} & UART0 & TX 43, RX 44 & разъём «UART» на плате \\
\code{Serial1} & UART1 & задаются в \code{begin()} & свободен \\
\code{Serial2} & UART2 & задаются в \code{begin()} & свободен \\
\bottomrule
\end{tabular}
\end{center}

\begin{figure}[H]
\centering
\begin{tikzpicture}[node distance=10mm]
  \node[blk,fill=blkmem,minimum width=30mm,minimum height=12mm] (pc) {Компьютер\\ \scriptsize Serial Monitor};
  \node[blk,fill=blkcpu,minimum width=30mm,minimum height=12mm,right=40mm of pc] (esp) {ESP32-S3\\ \scriptsize \code{Serial}};
  \draw[arr,accent] ($(pc.east)+(0,0.25)$) -- node[above,font=\ttfamily\scriptsize]{"on\textbackslash n"} ($(esp.west)+(0,0.25)$);
  \draw[arr,okgreen] ($(esp.west)+(0,-0.25)$) -- node[below,font=\ttfamily\scriptsize]{"OK\textbackslash n"} ($(pc.east)+(0,-0.25)$);
  \node[font=\sffamily\scriptsize,below=1mm of pc] {\code{Serial.readStringUntil()}};
  \node[font=\sffamily\scriptsize,below=1mm of esp] {\code{Serial.print()}, \code{printf()}};
\end{tikzpicture}
\caption{Обмен строками между компьютером и платой}
\end{figure}

Основные функции:
\begin{itemize}
  \item \code{Serial.begin(115200)} --- открыть порт; скорость в Serial Monitor должна совпадать.
  \item \code{Serial.print()}, \code{println()}, \code{printf()} --- отправить текст.
  \item \code{Serial.available()} --- сколько байт пришло и ждёт чтения.
  \item \code{Serial.readStringUntil('\textbackslash n')} --- прочитать строку до символа перевода строки.
\end{itemize}

\begin{tip}
В Serial Monitor выберите окончание строки \textsf{New Line}: тогда каждая отправленная команда будет
заканчиваться символом \code{\textbackslash n}, и \code{readStringUntil()} сразу поймёт, где конец команды.
\end{tip}

\subsection{Шаг проекта 2: команды с компьютера}

\progress{2}

Добавляем к шагу 1 обработку команд. Функция \code{handleSerial()} ничего не ждёт: если данных нет,
она сразу возвращается, и \code{loop()} продолжает крутить «пульс».

\codecap{Умная лампа, шаг 2 (полный скетч)}
\codeinput{lamp_step2}

\subsection{Второй UART: связь с другими устройствами}

Тот же протокол используют GPS-модули, модемы и другие микроконтроллеры. Благодаря матрице GPIO
UART1 и UART2 можно вывести почти на любые выводы.

\begin{figure}[H]
\centering
\begin{tikzpicture}
  \node[blk,fill=blkcpu,minimum width=3cm,minimum height=2.6cm] (a) {ESP32-S3};
  \node[blk,fill=blkper,minimum width=3cm,minimum height=2.6cm,right=5cm of a] (b) {GPS / другой МК\\ (3,3 В логика)};
  \draw[arr] ($(a.east)+(0,0.7)$) node[above right,font=\scriptsize]{TX (\pin{17})} -- ($(b.west)+(0,-0.1)$) node[below left,font=\scriptsize]{RX};
  \draw[arr] ($(b.west)+(0,0.7)$) node[above left,font=\scriptsize]{TX} -- ($(a.east)+(0,-0.1)$) node[below right,font=\scriptsize]{RX (\pin{18})};
  \draw[thick] ($(a.east)+(0,-0.8)$) node[below right,font=\scriptsize]{GND} -- ($(b.west)+(0,-0.8)$) node[below left,font=\scriptsize]{GND};
\end{tikzpicture}
\caption{TX одного устройства соединяется с RX другого; земли объединяются}
\end{figure}

\codecap{Мост: всё, что пришло с модуля, пересылается на компьютер, и наоборот}
\codeinput{l2_bridge}

\begin{caution}
Модули с 5-вольтовой логикой (часть Arduino, некоторые GPS) нельзя подключать к RX ESP32-S3 напрямую ---
используйте делитель или преобразователь уровней.
\end{caution}

\begin{summary}
\begin{itemize}[leftmargin=*]
  \item UART передаёт байты кадрами: старт-бит, 8 бит данных, стоп-бит; скорости сторон должны совпадать.
  \item \code{Serial} --- главный инструмент отладки: будем пользоваться им во всех следующих уроках.
  \item Команды читаем неблокирующе: проверяем \code{available()} и только тогда читаем строку.
\end{itemize}
\end{summary}

\begin{tasks}
\begin{enumerate}
  \item Добавьте в лампу команду \code{blink N}: мигнуть лампой N раз (без \code{delay()}, по образцу урока 1).
  \item Выводите раз в 10 секунд строку с временем работы, не мешая приёму команд.
  \item Соедините две платы по UART и передавайте команды лампе с другой платы.
\end{enumerate}
\end{tasks}
